\paragraph{Calculation and Testing Sequence}
Scheme~\ref{scheme:simulation-workflow} summarizes the following sequence for each parameter case.
\begin{enumerate}
  \item \textbf{Construct the domain and mesh.} Select the case from Table~\ref{tab:geom_cases}, varying only the nominated angle or height while retaining the companion parameters. Construct the capsule geometry, or retain the dimple array and change the flow-entry direction for its angle sweep. Generate the fluid-domain poly-hexacore mesh with near-wall prism layers and the mesh-quality controls in Section~\ref{sec:mesh}. Identify the inlet, outlet, two heated walls and lateral periodic pair. Evaluate $A_0$, $S_1$, $S_2$ and $D_e$, then calculate the reference speed for $\mathrm{Re}=4282$.
  \item \textbf{Assign models and solve a period.} Apply the constant water properties, energy equation, SST $k$--$\omega$ model, coupled pressure--velocity scheme and second-order discretization listed in Table~\ref{tab:bc-settings}. Starting from the current inlet data, iterate the steady equations and check the configured residual targets. Velocity, mass-flow, pressure-drop and thermal-report trends provide separate checks on physical-output stability; low residuals alone are insufficient.
  \item \textbf{Transfer the outlet profiles.} At the end of one period, export the outlet velocity components, $k$ and $\omega$, map them to the inlet, and apply the saved profile-scaling expressions before the next period. The velocity field is scaled to the reference mean speed and its in-plane direction is corrected to the prescribed flow-entry direction. The velocity scale factor is denoted as $a$, the inlet velocity component $u_x$, $u_y$, $u_z$, and $\omega$ will be scaled by $a$, while $k$ will be scaled by $a^2$. Keep $T_{in}=293$~K and the wall temperatures unchanged. Compare available profile and integral-output records between periods to assess hydrodynamic repeatability.
  \item \textbf{Reduce the final fields consistently.} Batch post-processing extracts the mesh areas, area-averaged inlet/outlet pressures, separate wall heat rates and fluid-domain mass-average static temperature from the saved case/data pairs. The same expressions in Section~\ref{sec:numerical_methods} give $D_e$, $u$, $h_{ave}$, Nu and $f_F$; Equations~\eqref{eq:PEC} and~\eqref{eq:red} give PEC and RED. Signed mass and enthalpy fluxes supply separate conservation checks. Store outputs by case label, retain the geometric and flow-direction identifiers, and compare the four sweeps and the targeted combined dimple case. This automated data-reduction step does not imply that the original geometry generation and solver loop were fully automated.
\end{enumerate}

\begin{scheme}[H]
\centering
\begin{tikzpicture}[
  font=\small,
  process/.style={draw=black!65,rounded corners=2pt,fill=black!3,
    text width=6.9cm,align=center,inner sep=6pt,minimum height=1.0cm},
  check/.style={diamond,draw=black!65,fill=black!3,aspect=3.2,
    text width=4.0cm,align=center,inner sep=1pt},
  transfer/.style={draw=black!65,rounded corners=2pt,fill=blue!5,
    text width=3.0cm,align=center,inner sep=6pt},
  arrow/.style={-{Stealth[length=2.2mm]},semithick,draw=black!75},
  edgeword/.style={font=\footnotesize,fill=white,inner sep=2pt}
]
  \node[process] (prepare) at (0,0)
    {Select parameter case; construct fluid domain\\
     Mesh with prism layers; check zones and quality};
  \node[process,below=5mm of prepare] (setup)
    {Set $D_e$ and $u$ for $\mathrm{Re}=4282$; apply models\\
     Lateral periodic pair; $T_{in}=293$~K, $T_w=333$~K};
  \node[process,below=5mm of setup] (solve)
    {Steady solver iterations within one period $n$\\
     Use current inlet velocity, $k$ and $\omega$ profiles};
  \node[check,below=5mm of solve] (residual)
    {Configured residual\\targets met?};
  \node[process,below=5mm of residual] (monitor)
    {Inspect $h_{ave}$, $\Delta p$, $\Delta k$,and $\Delta \omega$ \\
     *$\Delta p$, $\Delta k$,and $\Delta \omega$ between inlet and outlet};
  \node[check,below=5mm of monitor] (blocks)
    {Are all the monitored parameters monotonic or oscillatory converged?};
  \node[transfer,right=9mm of blocks] (map)
    {Export profiles\\from outlet\\
     Map and scale\\$\mathbf{u}$, $k$, $\omega$ at inlet\\
     Keep $T_{in}=293$~K\\
     Next Period};
  \node[process,below=5mm of blocks] (reduce)
    {Save final fields and monitor records\\
     Batch reduction: $Q_1,Q_2,T_{f,ave},h_{ave},\mathrm{Nu},f_F$};
  \node[process,below=5mm of reduce] (compare)
    {Repeat for remaining parameter cases\\
     Compare Nu, $f_F$, PEC and RED on the same basis};

  \draw[arrow] (prepare) -- (setup);
  \draw[arrow] (setup) -- (solve);
  \draw[arrow] (solve) -- (residual);
  \draw[arrow] (residual.west) -- ++(-8mm,0) |- (solve.west)
    node[edgeword,pos=0.22,left] {No: continue};
  \draw[arrow] (residual) -- node[edgeword,right] {Yes} (monitor);
  \draw[arrow] (monitor) -- (blocks);
  \draw[arrow] (blocks.east) -- node[edgeword,above] {Yes} (map.west);
  \draw[arrow] (map.north) |- (solve.east);
  \draw[arrow] (blocks) -- node[edgeword,right] {No} (reduce);
  \draw[arrow] (reduce) -- (compare);
\end{tikzpicture}
\caption{Computational workflow and streamwise quasi-periodic iteration loop. Steady solver iterations are nested inside the reported profile-transfer sequence; the lateral periodic pair remains active throughout. Temperature is not transferred from outlet to inlet. Physical-output stability is assessed separately from the period count. The final batch-reduction stage refers to the revised post-processing of the saved solutions.}
\label{scheme:simulation-workflow}
\end{scheme}
\clearpage
